%%%%%%%%%%%%%%%%%%%%%%%%%%%%%%%%%%%%%%%%%%%%%%%%%%%%%%%%%%%%%%%%%%%%%%%%%%%%
% forced-objects-pm.tex
%
% Philosophia Mathematica submission format: article class, author-year
% citations, footnotes, British spelling, an objections-and-replies section.
% Content is the same argument as ../forced-objects.tex (v5), re-addressed to a
% philosophy-of-mathematics readership.
%
% The honesty ledger travels with the paper:
%   * "checked" = Lean 4 declaration, no `sorry`, one named axiom (O1), and O1
%     is consumed by no theorem;
%   * "argued" = prose argument, not formalised;
%   * "asserted" = written down as an obligation;
%   * the standard-maps condition (C0) is not formalised and is flagged;
%   * Question 3 is specified, not implemented, and the first implementation
%     was refuted by the paper's own examples.
%%%%%%%%%%%%%%%%%%%%%%%%%%%%%%%%%%%%%%%%%%%%%%%%%%%%%%%%%%%%%%%%%%%%%%%%%%%%

\documentclass[12pt,a4paper]{article}

\usepackage[T1]{fontenc}
\usepackage[utf8]{inputenc}
\usepackage{lmodern}
\usepackage[margin=1in]{geometry}
\usepackage{amsmath,amssymb,amsthm}
\usepackage{booktabs}
\usepackage{array}
\usepackage[hidelinks]{hyperref}
\usepackage[authoryear,round]{natbib}
\usepackage{enumitem}
\usepackage{microtype}

\newtheorem{theorem}{Theorem}
\newtheorem{definition}[theorem]{Definition}
\newtheorem{obligation}[theorem]{Obligation}

\setlength{\parskip}{0.35em}
\setlength{\parindent}{0pt}

\title{Where the Ladder Must Stand\\
\large A test for forced mathematical objects, three verdicts, and one measurement that failed}

\author{Yuanjie Liu\\
\small Independent researcher\\
\small \texttt{yuanjieliu65@gmail.com}}

\date{}

\begin{document}

\maketitle

\begin{abstract}
\noindent
Mathematicians often say that an object \emph{had to} exist: you cannot subtract
in the natural numbers, so the integers appear; you cannot solve $x^{2}+1=0$ in
the reals, so the complex numbers appear.  The same objects are also called
inventions, and both descriptions are usually offered without a test.  This paper
proposes one, in the form of questions that can be answered yes or no, and reports
what happened when the test was turned on the author's own work.  Question~1: can
the demand be met on the same elements?  Question~2: among all extensions, is
there a smallest one, determined up to a unique relabelling?  Question~3: is the
result new, judged against a catalogue of known objects?  If Question~2 is
answered yes, the object is pinned down: any two smallest extensions are the same
up to a relabelling, and the relabelling is unique.  Three steps of the number
tower pass Questions~1 and~2; the step from the rationals to the reals does not,
and is recorded as a named obligation instead; a counting invariant proposed
inside the same programme is refused, because it is already definable on the same
elements.  Question~3 broke: the implementation compared the sizes of carriers,
and the integers, the rationals and the complex numbers all have the same size,
so the test refused everything.  The paper argues that the first two questions
are worth keeping, that the third needs to compare structures rather than
carriers, and that a test which refuses its author's candidates is worth more
than one which admits them all.
\end{abstract}

\noindent\textbf{Keywords:} mathematical discovery; minimal extension; initial
object; rigidity; number systems; formal verification; negative result.

\begin{quote}
\itshape
``Whatever can be defined may be questioned.''
\par
\vspace{0.4em}
--- \emph{Galactic Empire} (Isaac Asimov)
\end{quote}

\section{The puzzle}

\subsection{A sentence with no test}

When mathematicians introduce a number system they usually say something causal:
\emph{we could not subtract, so we made the integers}.  The sentence is not a
description of the construction.  It is a claim about necessity, and it is the
claim that does the philosophical work.  It is also the claim that is easiest to
abuse.

Consider a program that announces a new number $\alpha$ with $\alpha^{2}=7$.
Nothing in the announcement is false.  The rationals contain no square root of
seven, so $\alpha$ does extend them; but the extension is the field everyone
already knows, with $\alpha$ written in place of $\sqrt{7}$.  A test that asks only
`was there a solution before?' will certify this as new, and will certify a great
deal else besides.

The same weakness runs through the standard justifications.  `The natural numbers
cannot solve $x+1=0$' is true.  It is also true of the one-element set, of a free
group, and of the field of rational functions.  Each of them contains a solution;
none of them is the integers.  So the sentence \emph{there was no solution,
therefore we added one} does not pick out an object.  It picks out a crowd.

\subsection{The literature settles what, not why}

The tower itself is settled mathematics.  The integers can be built as pairs of
naturals modulo a difference, or as the group completion of the additive monoid;
the rationals as ordered pairs of integers, or as the fraction field; the reals as
Dedekind cuts \citep{dedekind} or as equivalence classes of Cauchy sequences; the
complex numbers as pairs of reals \citep[following Hamilton, 1837]{ebb} or as the
quotient ring $\mathbb{R}[x]/(x^{2}+1)$.  Any two constructions give isomorphic
objects; \citet{landau} builds the tower from Peano's axioms with the standard
uniqueness statements, and \citet{ebb} survey several constructions of each system
together with their equivalence.

That literature settles \emph{what} the objects are, and it leaves the causal
sentence untouched.  Mathematicians do have a name for `had to exist' --- the
object is \emph{initial}, or \emph{free}, or \emph{universal}, in the relevant
category --- and the standard fact that an initial object is determined up to a
unique isomorphism is in \citet{MacLane1971}.  Model theory supplies a parallel
notion: when a model completion exists it is unique up to equivalence, and the
comparison embeddings are elementary and unique over the base model
\citep{EklofSabbagh1971,Hodges1993}.  So the ingredients are old, and this paper claims no new
mathematics for them.

Programs that propose mathematics have a long history, from Lenat's \textsc{AM}
\citep{lenat1976} to recent systems that use search or learning to generate
conjectures and constructions \citep{davies2021,raayoni2021,romera2024}.  Their
pattern is consistent: generate candidates, then decide which ones are worth
keeping.  The decision is made by a human, by a heuristic of interestingness, or by
a checker that tests truth or performance rather than novelty.  This paper is about
the decision step, and should be read as complementary to those systems: a
generator supplies candidates, a criterion decides which of them may be called
forced.

What is missing, as far as we can tell from the literature cited here, is a
\emph{test}: a short list of questions that a reader can put to a claim of the
form `this object had to exist' --- including claims made by a program ---
together with an honest account of which questions can be settled mechanically
and which cannot.  The textbook statement is about a fixed category.  A test has
to survive a change of category, because that is where the abuses live.

\subsection{What this paper does}

The paper builds such a test, states what it can and cannot settle, and then runs
it on the programme that produced it.  Three results are negative and are
reported as such: one step of the tower is not certified, one candidate is
refused, and the third question is not implemented.  The positive content is
small and is stated in the same breath: two questions work, they are
construction-invariant, and they are enough to refuse a candidate that a weaker
test would have published.

The order of the paper follows the argument.  \S\ref{sec:senses} distinguishes
four things that `forced' can mean and says which one is tested.
\S\ref{sec:test} states the test.  \S\ref{sec:theorem} explains why the second
question settles the object.  \S\ref{sec:verdicts} reports the three verdicts.
\S\ref{sec:failure} reports the measurement that failed.
\S\ref{sec:objections} answers six objections.  \S\ref{sec:conclusion} concludes.

\section{Four senses of `forced'}
\label{sec:senses}

The word carries at least four different claims, and a paper that does not say
which one it means can be attacked from any direction.

\begin{enumerate}[label=(\roman*),leftmargin=2.2em]
\item \emph{Metaphysical}: the object could not have failed to exist.
\item \emph{Conceptual}: the object is already contained in the concept we
      started from.
\item \emph{Practical}: no working alternative exists, given the state of the
      field.
\item \emph{Determined}: every acceptable answer is the same object up to a
      unique relabelling.
\end{enumerate}

This paper tests (iv) and nothing else.  (i) is not touched: we have no argument
that the integers could not have failed to exist, and we doubt that a formal test
could supply one.  (ii) is not touched either; our demands are data, not
concepts, and a demand can be met in a way that tells us nothing about the
concept that motivated it.  (iii) is a claim about a community, and it changes
when the community changes; it is not the kind of thing a proof can settle.

The fourth sense is weaker than the first three, and it is the only one that can
be settled by reading a proof.  It is also the one that Benacerraf's argument
leaves standing.  \citet{benacerraf1965} observed that if numbers are sets, then the
natural number $4$ could be the von Neumann $4$ or the Zermelo $4$, and nothing in
arithmetic tells us which; so either the identification is arbitrary or numbers
are not objects at all.  Structuralism replies that numbers are positions in a
structure rather than objects with an inner nature \citep{shapiro1997}.  If that is
right, then `the same position' is exactly determination up to relabelling, and
nothing stronger is on offer.

We therefore do not claim to have shown that the integers had to exist.  We claim
to have shown that, given the demand and the company it keeps, nothing else could
have answered it.  That is not a hedge; it changes what counts as a refutation.  A
reader who wants metaphysical necessity will find the test silent, and should.  A
reader who wants to know whether a program's `discovery' is a renaming will find
it decisive.

\section{The test}
\label{sec:test}

\subsection{Three pieces of data}

Before any question can be asked, three things must be fixed.

\begin{enumerate}[label=(\alph*),leftmargin=2.2em]
\item An \emph{ambient theory}: what the new object is allowed to carry.
\item A \emph{reading} of the old structure inside the new one.
\item A \emph{demand}: a constant to add, a property to impose, or both.
\end{enumerate}

`Minimal extension' is not an absolute notion, and the reason is easy to see.
Given only a system and a missing solution, one may always take a trivial
extension (throw in a formal symbol and nothing else) or a chaotic one (throw in a
proper class of symbols), and there is no order between them that would single
out a winner.  A minimum becomes definable only once (a)--(c) are fixed.

The first step of the tower shows why (a) matters, and it is worth stating slowly.
The demand is `$x+1=0$ must become solvable'.  If we are allowed to stay among
structures where addition is only a semigroup operation, we may add a formal
$\varepsilon$ with $\varepsilon+1=0$ and stop.  The result solves the equation and
is not the integers: its addition is not a group operation.  If instead we first
agree to work among commutative rings, the same demand produces the integers, and
produces nothing else.  The difference is not in the demand; it is in the company
the demand is allowed to keep.  Changing the ambient theory changes the object,
and the change of theory is the real content of the step.

\subsection{Question 1: can the demand be met on the same elements?}

The question is about the elements, not about the solution.  The wording is doing
work.  Both the natural numbers and the integers are countable, so any ring
structure on the integers can be carried over to the natural numbers by a
bijection.  The statement `the natural numbers admit no ring structure' is
therefore \emph{false}, and an early version of our own formalisation asserted it.
The true statement keeps the old structure fixed: no ring structure on
$\mathbb{N}$ whose addition and multiplication extend the usual ones solves
$x+1=0$.  The episode is worth recording because the test caught its author.

A gap, then, is not `this carrier admits no such structure'.  It is `this carrier,
with the structure it already has, cannot support the required one'.  The
distinction is the formal counterpart of a platitude: you cannot solve an
equation by changing what the symbols mean.

\subsection{Question 2: is the extension the smallest one?}

Among all extensions that meet the demand, is there one that maps into every
other, and does exactly one such map exist?  When the answer is yes we call the
extension \emph{minimal}.  The name is traditional --- `the minimal field
containing a root' --- but the property is not smallness; it is universality, and
universality is what makes the object determined.  The next section states what
follows from it.

Two clarifications, because `smallest' invites misreading.  First, the comparison
is between \emph{extensions}, not between carriers: a smaller carrier is neither
necessary nor sufficient.  Second, `smallest' is relative to the triple fixed
above; the same system and the same demand can have different smallest extensions
under different ambient theories, and that is a finding rather than a defect.

\subsection{Question 3: is it new?}

New relative to what is already known, not relative to the system's own state,
which can always be arranged to report novelty.  In practice this means an
explicit catalogue of known objects, supplied from outside.  The requirement is
not stylistic.  Whether an object is new is a fact about what is already known,
and a system that computes novelty from its own state will always find itself
original.

\S\ref{sec:failure} reports what happened when this question was implemented.  The
short version is that the implementation was wrong in a way that only the
examples revealed, and that the test as a whole is therefore a two-question test
with a third question specified and unimplemented.  We prefer to say that at the
outset rather than let a reader discover it.

\section{Why the second question settles the object}
\label{sec:theorem}

\subsection{The statement}

If an extension is smallest in the sense above, then:

\begin{enumerate}[label=(\arabic*),leftmargin=2.2em]
\item it maps into every other extension, and the map is unique;
\item any two smallest extensions are the same up to a relabelling of their
      elements, and that relabelling is itself unique;
\item the smallest extension has no symmetries of its own beyond the identity.
\end{enumerate}

None of this mentions rings, order, or completeness.  The three statements use
four properties of the surrounding category --- an identity map, a composition, an
underlying-function map, and the fact that maps agreeing on every value are equal
--- and nothing else.  That is why one argument covers the integers, the
rationals and the complex numbers at once, and why it can be checked once rather
than three times.

\subsection{What the theorem is worth}

We should be plain about this, because the alternative is to let a reader
discover it.  Once you \emph{define} `smallest' as `maps into everything,
uniquely', the first and third statements are restatements of the definition, and
the proof is four lines.  The content is not the proof.  The content is that this
is the right definition to have chosen: it is the only one on the table that does
not depend on how the object is constructed.  Two people who build the complex
numbers differently --- as pairs of reals, as a quotient ring, as a matrix algebra
--- get the same answer to Question~2, and their constructions are related by a
relabelling that is itself unique.  That is what it means for the object, rather
than the construction, to be determined.

A philosophical reader may object that this is precisely the structuralist
commonplace, and that a four-line proof of a commonplace is not worth a paper.  The
objection is right about the proof and wrong about the use.  What the four lines
buy is \emph{transfer}: because the statement is parameterised by the ambient
theory and the reading, the same receipt covers every step of the tower, and the
receipt can be checked by machine.  The interest lies in what the receipt refuses,
not in how hard it was to obtain.

\subsection{Pinning the witness}

For the step from the reals to the complex numbers, the demand is `add a constant
$j$ with $j\cdot j=-1$'.  The extension $(\mathbb{C},i)$ is smallest, and it has
no symmetry beyond the identity.  Now drop the constant from the demand and ask
only that \emph{some} square root of $-1$ exists.  The same field is still an
extension, but it is no longer smallest: the comparison is no longer unique.
Identity and complex conjugation are two distinct maps over $\mathbb{R}$, and both
are allowed; pinning $i$ is what removes the second one.

The moral is general.  Finding a solution and pinning it down are different
achievements, and only the second produces an object.  A test that accepts the
unpinned version will accept a field with a symmetry group, which is exactly the
ambiguity it was supposed to remove.

\subsection{What the theorem does not settle}

The demand `the order is complete' behaves differently from `$j^{2}=-1$', and the
difference is worth stating in words rather than symbols.  An equation like
$j^{2}=-1$ is preserved by every ring map: if $j$ satisfies it, so does the image
of $j$.  A condition like `every bounded non-empty set has a least upper bound' is
not preserved by ring maps, and it has no term model: there is no expression you
can write down, quotient by a relation, and be done.  What follows is that the
smallest extension cannot be \emph{computed} the way the complex numbers can.  It
does not follow that no smallest extension exists.

An earlier draft of this paper drew the stronger conclusion.  It was wrong, and
the correction matters for how the fourth step of the tower is reported: the
asymmetry is about computation and about proof debt, not about a metaphysical
difference between adjoining a root and completing an order.

\section{Three verdicts}
\label{sec:verdicts}

We ran the test on four steps of the number tower and on two further pieces of the
same programme.  Table~\ref{tab:verdicts} gives the summary; the text explains the
interesting rows.

\begin{table}[t]
\centering
\small
\begin{tabular}{@{}lccc@{}}
\toprule
Case & Same elements & Smallest & New \\
\midrule
$\mathbb{N}\to\mathbb{Z}$ & proved & proved & not used \\
$\mathbb{Z}\to\mathbb{Q}$ & proved & proved & not used \\
$\mathbb{Q}\to\mathbb{R}$ & proved & \emph{asserted} & not used \\
$\mathbb{R}\to\mathbb{C}$ & proved & proved & not used \\
counting invariant & \emph{fails} (argued) & --- & --- \\
\bottomrule
\end{tabular}
\caption{What passed, what did not, and what could not be judged.  `Proved' means
machine-checked modulo the standard-maps condition discussed in
\S\ref{sec:objections}; `asserted' means written down as an obligation and not
proved; `argued' means a prose argument, not a formalised verdict.  No entry in
the `New' column was produced by an implementation.}
\label{tab:verdicts}
\end{table}

\subsection{Three steps that pass}

The integers are the smallest extension of the naturals once we agree to work
among commutative rings; the rationals are the smallest extension of the integers
once we agree to work among fields; and $(\mathbb{C},i)$ is the smallest extension
of the reals once we pin the square root of $-1$.  The underlying facts are
classical \citep{ebb}; the formulation as a smallest extension, with the old
structure pinned, is our packaging.  What the test adds is that the three are
instances of one statement, and that in each case the answer does not depend on
the construction.  For the integers, the side condition that the extension really
contains the old system also gives characteristic zero for free, which is a
pleasant sign that the definition is carrying its weight.

\subsection{One step we could not certify}

For the step from the rationals to the reals we prove the failure --- the rationals
are not complete --- and we prove that every real is generated by a cut of
rationals below it.  We do not prove that the reals are the smallest complete
ordered field over the rationals.  That statement is written down as a named
obligation in the formal development.

We think the obligation is true and provable; a page of standard argument about
complete ordered fields should do it.  We did not do it.  Nothing in the
development consumes the obligation, so nothing collapses if it is later proved or
refuted; what would change is the status of the fourth row of
Table~\ref{tab:verdicts}, which is exactly what should change.

\subsection{One candidate we refused}

The programme also proposed an invariant of integers, meant to capture the
`tension' between two numbers: count the primes dividing one but not the other,
subtract the count in the opposite direction, and call the result $\tau(a,b)$.  The
proposal was that the surface $\tau=0$ organises arithmetic, and that $\tau$ is a
new object in its own right.

The axiomatic part worked out cleanly.  Under four natural assumptions the
invariant is unique up to sign and equals a difference of counts,
\begin{equation}
\tau(a,b) = \omega(a) - \omega(b),
\end{equation}
where $\omega(n)$ is the number of distinct prime factors of $n$.  That is the
whole content: $\tau$ knows how many prime factors each side has and never which
ones.  Since $\omega$ is a century-old function and the formula is definable on the
same elements, the test returns \emph{refused} at Question~1.  There is no new
object here; there is a theorem about $\omega$, and a clear statement of why the
construction cannot say more.

Two qualifications.  The refusal is argued by hand: the counting part of the
development contains no demand and no catalogue, so `Question~1 returns refused'
is a mathematical observation, not a machine-checked verdict.  And the refusal
does not go through Question~3: the first draft said the novelty score was zero
there, which neither follows from Question~1 nor is available, since no catalogue
exists.

The episode is still the most useful thing in the paper.  It is the case where the
test said no to us, and it came with a reason: the invariant collapses to a count,
and counts are already on the same elements.  A weaker test would have produced a
paper claiming a discovery.

\section{The measurement that failed}
\label{sec:failure}

Question~3 is the one we could not deliver, and the failure is instructive enough
to describe in detail.

\subsection{Two scores}

The design had two scores.  One was a weighted average of five qualities ---
coverage, relatedness, closure gain, minimality, novelty --- meant to \emph{rank}
candidates.  The other was a product of four --- novelty, consistency,
generativity, compression --- meant to \emph{admit} them.  The distinction is real:
a weighted average gives a high mark to a candidate whose novelty is zero but whose
other qualities are excellent, which is precisely the $\sqrt{7}$ story; a product
gives it zero.

The claim that the two scores are not interchangeable is machine-checked in the
only form available: there is a score with novelty zero and generativity $4/5$,
and a snapshot with the same entries whose product is zero.  We stress what this
does and does not show.  It shows that the two formulas are not interchangeable.
It does not model renaming at all: no catalogue, no candidate, and no map between
objects appears in those two statements.  The $\sqrt{7}$ story remains an
argument, not a theorem.

\subsection{Why the implementation failed}

Modelling renaming is where the implementation failed.  To decide whether a
candidate was new, the code compared the \emph{sizes} of carriers: if a
candidate's carrier had the same cardinality as a known object's carrier, it was
called a renaming.  The integers, the rationals and the complex numbers are all
countably infinite.  One countable object in the catalogue therefore makes all
three renamings of it, and the gate assigns all three a novelty of zero.  The gate
refuses exactly the objects the rest of the paper certifies.

We did not patch this.  Comparing sizes can only produce false refusals, never
false admissions, so the risk runs in the opposite direction from the one our
first draft described.  A correct comparison is between structures, not carriers,
which requires the catalogue to carry its own maps; that is a specification, and
we have not implemented it.

\subsection{What follows}

Three consequences, stated plainly.  First, no admission decision in this paper
uses Question~3, and the trace distributed with the code prints novelty values
that were entered by hand.  Second, the product structure and the zero-factor
lemmas survive, because they are statements about a product; what fails is the
\emph{measurement} of one factor.  Third, the failure is evidence for the test
rather than against it: the example that was supposed to demonstrate the gate
demolished it, and the demolition was visible only because the gate had been
applied to objects whose status we already knew.

\section{Objections and replies}
\label{sec:objections}

\subsection{`Your verdicts are relative to chosen data, so they are arbitrary'}

Two parts of the objection should be separated.  The first is true and is stated
in the paper: every verdict is relative to the ambient theory, the reading, and
(for Question~3) the catalogue.  Change any of them and the verdict can change.
The second part --- that relativity makes the verdict arbitrary --- does not
follow.  The data are published, so a reader can recompute the verdict with
different data and see whether it changes.  A relative verdict that anyone can
recheck is a different thing from an arbitrary one.

There is also a point of substance behind the relativity.  Changing the ambient
theory changes the object: the same demand for subtraction produces the integers
among rings and a semiring monster among semirings.  That is not a defect of the
test; it is the phenomenon the test was built to expose.  A test that gave the
same answer under every ambient theory would be insensitive to precisely the
choice that does the work.

\subsection{`A degenerate framework makes every object smallest'}

This is the most serious objection, and the paper concedes it.  The four
properties that the main argument uses do not say that the maps between
structures are the structure-preserving ones.  A degenerate choice --- one map
between any two structures --- satisfies all four, and in that setting every
object is both smallest and largest, so the test says nothing about anything.

The paper therefore states a precondition, called C0: the framework is the
standard one of the theory at hand.  It is not provable inside the framework,
because `being the standard maps' is a schema rather than a sentence; it is
checkable in each instance, and in our instances it holds by construction, since
the maps are defined to be the structure-preserving ones.  Every statement that
needs C0 is marked as such.

A test with a stated precondition is better than a test with an unstated one.  But
the objection does have a residue: the general theorem is less general than its
statement suggests, because its content lives in the instances rather than in the
abstraction.  We accept that description.

\subsection{`Benacerraf already showed that determination is all you get'}

Correct, and the paper says so.  If numbers are positions in a structure, then
determination up to a unique relabelling is the most that could be on offer, and
the test is not adding a metaphysical result.  What it adds is a procedure: given
a candidate and a catalogue, a reader can decide whether the word `discovery' is
warranted, and can be shown the step where the decision is made.

Whether that is philosophy is a fair question.  Our answer is that it is the part
of the philosophy that can be done with a proof rather than a intuition, and that
the two refusals reported here --- one by proof, one by argument --- are the kind
of result the procedure produces.  A reader who wants the metaphysics will not find
it, and we have tried not to pretend otherwise.

\subsection{`Your main theorem is a definition'}

Yes.  The paper says so in the section that states it, and repeats it in the
conclusion.  The theorem is worth stating because it is the receipt that the
definition is well posed, and because the receipt transfers across theories.  It is
not worth stating as a discovery, and we do not.

The objection would bite if the paper's contribution were the theorem.  It is not:
the contribution is the triple of questions, the two refusals, and the report of
the failure.  A reader who judges the paper by the length of its main proof will
judge it correctly and will not want to read it.

\subsection{`One refusal of your own candidate is not evidence'}

It is not much evidence, and the paper says that too.  A single refusal, argued by
hand, of a candidate the author was already prepared to abandon, is close to the
weakest form of evidence a test can produce.  What would count as stronger: a
candidate proposed by someone else and believed by them, refused by the test; an
independent user applying the test to a case neither of us chose; or a repaired,
structure-comparing gate that still refuses an object the test ought to admit.

None of these exists yet.  We list them because a reader is entitled to know what
would change our mind, and because the list is short enough to be a research
programme rather than a wish.

\subsection{`Machine-checking here is decoration'}

Some of it is, and the paper tries to say which parts.  The development is in
Lean~4 \citep{lean} with Mathlib \citep{mathlib}.  The main theorem's proof is
short enough that checking it adds little; the four failures are textbook; the
two-score separation is arithmetic.  Three things were not decoration.  The proof
assistant forced the definition of `gap' to be stated in a form that is not false,
which is how the ring-structure error was found.  It made the four-line nature of
the main argument visible, so that the paper could not pretend otherwise.  And it
turned a claim about two scoring formulas into a statement a reader can check in a
minute.  A tool whose main contribution is to make the author's mistakes visible is
worth its cost in a paper of this kind.

\section{Conclusion}
\label{sec:conclusion}

The practical lesson is about the order in which to explain an object.  The usual
order is: here is the object, here are the properties that make it necessary.  The
order that survived the test is the reverse.  First the failure that cannot be
repaired in place.  Then the company the repair is allowed to keep, because the
same failure repaired in different company gives different objects.  Then the
smallest repair, which is determined once the company is fixed.  Whatever is left
over after those three steps is not forced, and should be called a construction.

The second lesson is that a negative result is cheap and worth publishing.  The
counting invariant was, for a while, the centrepiece of a different paper.  Running
it through Question~1 took an afternoon and produced a theorem about $\omega$
instead.  That is a better outcome than a paper claiming a discovery, and it cost
less.

The third lesson is about what a test can be.  This one is two questions and a
failure.  It does not generate candidates, it does not decide its own questions,
and it cannot certify the fourth step of the tower.  What it can do is stop a
particular kind of sentence --- `there was no solution, therefore I discovered an
object' --- from being uttered without a receipt.  On the evidence of this paper,
that is enough to refuse one candidate and to expose one measurement, and not
enough to claim more.

\appendix

\section{What is checked, argued, and assumed}

\begin{table}[h]
\centering
\small
\begin{tabular}{@{}p{0.34\linewidth}p{0.14\linewidth}p{0.42\linewidth}@{}}
\toprule
Claim & Status & Note \\
\midrule
Uniqueness and rigidity of minimal extensions & checked & modulo C0 \\
Four failures of the tower & checked & natural, integer, rational, real \\
$\mathbb{Z}$, $\mathbb{Q}$, $(\mathbb{C},i)$ are minimal & checked & modulo C0 \\
Characteristic zero of the field extension & checked & follows from injectivity \\
Two-score separation & checked & arithmetic, not a model of renaming \\
Counting collapse to $\omega(a)-\omega(b)$ & checked & four axioms \\
Refusal of the counting invariant & argued & no demand, no catalogue instance \\
Real is the smallest complete ordered field over $\mathbb{Q}$ & asserted & obligation, consumed by nothing \\
Preservation of equations vs closure conditions & argued & about the language, not inside it \\
Novelty gate & refuted & carrier counting; repair specified, not built \\
\bottomrule
\end{tabular}
\caption{The status of every substantive claim.  `Checked' means a Lean~4
declaration with no \texttt{sorry} and no named axiom beyond the obligation
above; `argued' means prose; `asserted' means written down and not proved;
`refuted' means the implementation was shown wrong by the paper's own examples.}
\end{table}

\bibliographystyle{plainnat}
\bibliography{../nods}

\end{document}
